\chapter{Formül Özeti ve Hızlı Başvuru}
\label{app:formulas}
Bu ek, kitabın bölümlerinde doğrudan kullanılan temel formülleri hızlı
başvuru amacıyla bir araya getirir. Kritik değer tabloları
Ek~\ref{app:distribution-tables}'te verilmiştir.


\section{Betimsel İstatistik ve Sağlam Ölçüler}
\begin{align*}
\bar{x} &= \frac{1}{n}\sum_{i=1}^{n}x_i, \\
s^2 &= \frac{1}{n-1}\sum_{i=1}^{n}(x_i-\bar{x})^2, \\
\IQR &= Q_3-Q_1, \\
\text{Aykırı değer adayları} &: \quad x<Q_1-1{,}5\IQR\; \text{veya}\; x>Q_3+1{,}5\IQR.
\end{align*}

\section{Beklenen Değer ve Örnekleme Değişkenliği}
\begin{align*}
\E[aX+b] &= a\E[X]+b, \\
\Var(X) &= \E[X^2]-\{\E[X]\}^2, \\
\Var(aX+b) &= a^2\Var(X), \\
\E[\bar X] &= \mu, \qquad \Var(\bar X)=\frac{\sigma^2}{n}, \\
\E[\hat p] &= p, \qquad \Var(\hat p)=\frac{p(1-p)}{n}.
\end{align*}

\section{İkili Gözlemler ve Örneklem Oranı}
\begin{align*}
X\sim\mathrm{Bernoulli}(p) &: \quad \P(X=1)=p, \quad \E[X]=p, \quad \Var(X)=p(1-p), \\
X\sim\mathrm{Binom}(n,p) &: \quad \P(X=k)=\binom{n}{k}p^k(1-p)^{n-k}, \\
&\quad \E[X]=np, \quad \Var(X)=np(1-p), \qquad \hat p=\frac{X}{n}.
\end{align*}

\section{Sık Kullanılan Güven Aralıkları}
\begin{align*}
\mu\; (\sigma \text{ biliniyor}) &: \quad \bar{x} \pm z_{1-\alpha/2}\frac{\sigma}{\sqrt{n}}, \\
\mu\; (\sigma \text{ bilinmiyor}) &: \quad \bar{x} \pm t_{1-\alpha/2,n-1}\frac{s}{\sqrt{n}}, \\
p &: \quad \hat{p} \pm z_{1-\alpha/2}\sqrt{\frac{\hat{p}(1-\hat{p})}{n}}.
\end{align*}

\section{Temel Test İstatistikleri}
\begin{align*}
T &= \frac{\bar{X}-\mu_0}{S/\sqrt{n}}, \qquad \nu=n-1, \\
T_{\text{Welch}} &= \frac{\bar X_1-\bar X_2}
{\sqrt{S_1^2/n_1+S_2^2/n_2}}, \\
T_{\text{eş}} &= \frac{\bar D}{S_D/\sqrt n}, \\
\chi^2 &= \sum_i \frac{(O_i-E_i)^2}{E_i},
\qquad E_{ij}=\frac{R_iC_j}{N}.
\end{align*}

\ifimoders\else
\section{Bootstrap ve Rastgeleleştirme}
\begin{align*}
\SE_{\mathrm{boot}}(\hat\theta)
&=\sqrt{\frac{1}{B-1}\sum_{b=1}^{B}
(\hat\theta_b^*-\bar\theta^*)^2}, \\
\text{Yüzdelik yüzde 95 GA}
&=[\hat\theta^*_{0{,}025};\hat\theta^*_{0{,}975}], \\
p_{\mathrm{rand}}
&=\frac{1+\sum_{b=1}^{B}\Ind{|T_b^*|\geq|T_{\text{göz}}|}}{B+1}.
\end{align*}
\fi

\section{Standart Hata ve Etki Büyüklükleri}
\begin{align*}
\SE(\bar X) &= \frac{\sigma}{\sqrt n}, \qquad
\widehat{\SE}(\bar X)=\frac{s}{\sqrt n}, \\
\widehat{\SE}(\hat p) &= \sqrt{\frac{\hat p(1-\hat p)}{n}}, \\
d_{\mathrm{tek}} &= \frac{\bar x-\mu_0}{s}, \\
d_z &= \frac{\bar d}{s_D}, \\
V &= \sqrt{\frac{\chi^2}{N\min(r-1,c-1)}}.
\end{align*}

\section{Korelasyon ve Basit Doğrusal Regresyon}
\begin{align*}
r &= \frac{\sum_i(x_i-\bar{x})(y_i-\bar{y})}
{\sqrt{\sum_i(x_i-\bar{x})^2\sum_i(y_i-\bar{y})^2}}, \\
b_1 &= \frac{\sum_i (x_i-\bar{x})(y_i-\bar{y})}{\sum_i (x_i-\bar{x})^2}, \\
b_0 &= \bar{y}-b_1\bar{x}, \\
\hat y_i &= b_0+b_1x_i, \qquad e_i=y_i-\hat y_i, \\
R^2 &= r^2 \qquad \text{(tek açıklayıcılı doğrusal modelde)}.
\end{align*}

\ifimoders\else
\section{Çoklu Doğrusal Regresyon}
\begin{align*}
Y_i &= \beta_0+\beta_1X_{1i}+\cdots+\beta_KX_{Ki}+\varepsilon_i, \\
\hat Y_i &= b_0+b_1X_{1i}+\cdots+b_KX_{Ki}, \\
e_i &= y_i-\hat y_i, \\
R^2 &= 1-\frac{\SSE}{\SST}, \\
R^2_{\mathrm{düz}}
&=1-(1-R^2)\frac{n-1}{n-K-1}, \\
\VIF_j &= \frac{1}{1-R_j^2}.
\end{align*}

İki değişkenli etkileşim modelinde
\[
Y=\beta_0+\beta_1X+\beta_2Z+\beta_3XZ+\varepsilon;
\]
$X$ eğimi $Z=0$ için $\beta_1$, $Z=1$ için $\beta_1+\beta_3$'tür.

\section{Tek Yönlü Varyans Analizi}
\begin{align*}
SS_{\mathrm{toplam}}
&=SS_{\mathrm{gruplar}}+SS_{\mathrm{hata}}, \\
SS_{\mathrm{gruplar}}
&=\sum_{j=1}^{k}n_j(\bar x_j-\bar x_{..})^2, \\
SS_{\mathrm{hata}}
&=\sum_{j=1}^{k}\sum_{i=1}^{n_j}(x_{ij}-\bar x_j)^2, \\
MS_{\mathrm{gruplar}}
&=\frac{SS_{\mathrm{gruplar}}}{k-1}, \qquad
MS_{\mathrm{hata}}=\frac{SS_{\mathrm{hata}}}{N-k}, \\
F&=\frac{MS_{\mathrm{gruplar}}}{MS_{\mathrm{hata}}}, \\
\eta^2&=\frac{SS_{\mathrm{gruplar}}}{SS_{\mathrm{toplam}}}, \\
\omega^2
&=\frac{SS_{\mathrm{gruplar}}-(k-1)MS_{\mathrm{hata}}}
{SS_{\mathrm{toplam}}+MS_{\mathrm{hata}}}.
\end{align*}
\fi